%% 学术研究汇报（Metropolis 现代深色主题）
%% 依赖 TeX Live 自带 beamertheme-metropolis
%% 编译: xelatex（推荐，中文）
\documentclass[UTF8,aspectratio=169,10pt]{ctexbeamer}
\usetheme{metropolis}
\usepackage{ctex}
\usepackage{graphicx}
\usepackage{amsmath,amssymb}
\usepackage{booktabs}
\usepackage{listings}

\setCJKmainfont{SimSun}
\setCJKsansfont{Microsoft YaHei}
\setCJKmonofont{FangSong}

\metroset{block=fill}
\definecolor{accent}{HTML}{2563EB}

\title{基于深度学习的轨迹预测方法}
\subtitle{组会 / 学术报告}
\author{姓名}
\institute{单位 / 实验室}
\date{\today}

\begin{document}

\maketitle

\begin{frame}{目录}
  \tableofcontents
\end{frame}

\section{研究背景}

\begin{frame}{研究动机}
  \begin{itemize}
    \item 自动驾驶对周围车辆轨迹预测的需求日益增长
    \item 传统模型难以捕捉交互与不确定性
    \item 深度学习提供了端到端建模能力
  \end{itemize}
\end{frame}

\begin{frame}{相关工作}
  \begin{columns}[T]
    \begin{column}{0.48\textwidth}
      \textbf{传统方法}
      \begin{itemize}
        \item 物理模型（IDM、MOBIL）
        \item 滤波与优化
      \end{itemize}
    \end{column}
    \begin{column}{0.48\textwidth}
      \textbf{数据驱动}
      \begin{itemize}
        \item RNN / LSTM
        \item Transformer
        \item 图神经网络
      \end{itemize}
    \end{column}
  \end{columns}
\end{frame}

\section{方法}

\begin{frame}{总体框架}
  输入历史轨迹 $\rightarrow$ 编码器 $\rightarrow$ 交互建模 $\rightarrow$ 解码器 $\rightarrow$ 多模态预测
\end{frame}

\begin{frame}{损失函数}
  \begin{align}
    \mathcal{L}
    &= \mathcal{L}_{\mathrm{reg}} + \lambda \mathcal{L}_{\mathrm{cls}} \\
    \mathcal{L}_{\mathrm{reg}}
    &= \frac{1}{N}\sum_{i=1}^{N}\left\lVert \hat{y}_i - y_i \right\rVert_2^2
  \end{align}
\end{frame}

\section{实验}

\begin{frame}{主要结果}
  \begin{table}
    \centering
    \begin{tabular}{lcc}
      \toprule
      方法 & ADE$\downarrow$ & FDE$\downarrow$ \\
      \midrule
      LSTM & 1.12 & 2.45 \\
      Transformer & 0.98 & 2.10 \\
      \textbf{Ours} & \textbf{0.91} & \textbf{1.96} \\
      \bottomrule
    \end{tabular}
  \end{table}
\end{frame}

\begin{frame}{小结与展望}
  \begin{itemize}
    \item 提出了交互感知的多模态轨迹预测框架
    \item 在公开数据集上取得竞争力结果
    \item 未来：可解释性、实时性、闭环评测
  \end{itemize}
\end{frame}

\begin{frame}[standout]
  谢谢！欢迎提问
\end{frame}

\end{document}
